\chapter{Zusammenfassung und Ausblick}
Der Schlussteil besteht aus einer kurzen Zusammenfassung und kann auch persönliche Bemerkungen, Bewertungen etc. enthalten. Zuerst wird somit das Vorher geschriebene knapp zusammengefasst. Anschließend werde die Ergebnisse reflektiert. Die Zusammenfassung setzt sich kritisch mit der Aufgabenstellung auseinander und beantwortet die Fragen:
\begin{itemize}
	\item Was sollte getan werden und was wurde tatsächlich erreicht?
	\item Wo gab es besondere Schwierigkeiten?
	\item Welche Teile der Aufgabenstellung konnten eventuell gar nicht bearbeitet werden und warum?
\end{itemize}

Der Ausblick kann Empfehlungen für eine weitere sinnvolle Fortsetzung des Projektes bzw. der Forschungsarbeit enthalten.
Solche Empfehlungen beruhen üblicherweise auf Erkenntnissen, die während der Bearbeitung des laufenden Projektes entstanden sind \cite[S.~44]{Hering2009}.

Es wird darauf hingewiesen, dass die wissenschaftliche Arbeit am Ende unbedingt von einer dritten Person korrekturgelesen werden sollte.
Dadurch können Fehlerarten wie falsche Wortwahl, Rechtschreibfehler. Zeichensetzungsfehler, Grammatikfehler, Stilfehler, inhaltliche Wiederholungen, Fehler im logischen Gedankenfluss, Layout-Fehler sowie sonstige Fehler aufgedeckt werden \cite[S.~174]{Hering2009}.
Zudem wird der Autor beim häufigen Lesen der eigenen Texte mit der Zeit "`betriebsblind"' gegenüber den eigenen Formulierungen.
Auch aus diesem Grund sollte von einem Außenstehenden überprüft werden, ob der Bericht für Projektfremde verständlich ist \cite[S.~130]{Hering2009}.